% Model training: decisions from the input of the models to the choice of the number of topics, each compared with
% Gordon (2026). Follows the preprocessing table (preprocessing_decisions.tex).
% Requires in the preamble:
%   \usepackage{pdflscape}  % landscape pages (rotated in the PDF viewer)
%   \usepackage{array}
%   \usepackage{longtable}  % the table runs over several pages

\begin{landscape}
{\small
% A section heading over all five columns
\newcommand{\steprow}[1]{\multicolumn{5}{|l|}{\textbf{#1}} \\}
\renewcommand{\arraystretch}{1.3}
\setlength{\LTcapwidth}{\linewidth}
\begin{longtable}{|>{\raggedright\arraybackslash}p{0.12\linewidth}
                  |>{\raggedright\arraybackslash}p{0.20\linewidth}
                  |>{\raggedright\arraybackslash}p{0.10\linewidth}
                  |>{\raggedright\arraybackslash}p{0.28\linewidth}
                  |>{\raggedright\arraybackslash}p{0.18\linewidth}|}
\caption{Model training: decisions from the input of the models to the number of topics, each compared with Gordon
(2026). \emph{Justification}: \emph{Motivated} = chosen for a stated reason; \emph{Gordon} = taken over from Gordon;
\emph{Default} = a library default or inherited setting that has not been tested; \emph{Open} = still to be decided.}
\label{tab:training} \\
\hline
\textbf{Decision} & \textbf{Setting} & \textbf{Justification} & \textbf{Comment} & \textbf{Gordon (2026)} \\
\hline
\endfirsthead
\hline
\textbf{Decision} & \textbf{Setting} & \textbf{Justification} & \textbf{Comment} & \textbf{Gordon (2026)} \\
\hline
\endhead
\multicolumn{5}{r}{\emph{continued on next page}} \\
\endfoot
\endlastfoot
%
% ---------------------------------------------------------------- input
\steprow{1 Input to the models}
\hline
Documents & 11,628 motions: the 13,943 extracted motions without 294 motions under 20 words, 642 later copies of
  duplicate texts and 1,379 near-copies ($\geq 80\%$ of three-word sequences shared with an earlier motion) & Motivated &
  The same motions in every model: the whole-corpus model and the four period models together. Median 160 words per
  motion after preprocessing (mean 329). & 14,973 chunks of about 1,000 words from 238 books. \\
\hline
Representation & Word counts per motion (bag of words); word order is ignored & Default & The standard input of LDA. &
  The same. \\
\hline
Vocabulary, whole-corpus model & Words in $\geq 10$ and $\leq 50\%$ of motions, no cap: 14,772 words, 91.8\% of the
  words of the final text & Motivated & The same rule as the period models, so the two designs differ only in the time split
  and the number of topics. & Words in $\geq 10$ and $\leq 50\%$ of documents, no cap. \\
\hline
Vocabulary, period models & Own vocabulary per period: words in $\geq 10$ and $\leq 50\%$ of the period's motions, no
  cap: 7,098 / 7,665 / 4,368 / 5,778 words & Gordon & Each period model is built on its own most relevant words.
  Gordon's documents are much longer than our motions. \textbf{Question:} is a minimum of 10 documents, calibrated on
  1,000-word chunks, reasonable for short motions? & The same rule, applied within each bin. \\
\hline
Upper limit & $\leq 50\%$ of documents & Open & In the whole corpus it removes 2 words, \textit{utbildning}
  (64\% of motions) and \textit{skola} (just over 50\%). \textbf{Question:} drop these central content words of an education
  corpus, or keep them? & The same limit; in long, varied chunks it mainly removes general words. \\
\hline
%
% ---------------------------------------------------------------- algorithm
\steprow{2 Algorithm and settings (the same for every model)}
\hline
Implementation & scikit-learn \texttt{LatentDirichletAllocation}, online variational Bayes & Motivated (open) &
  Chosen for practicality. The study is framed as an adaptation of Gordon's workflow, not a computational replication;
  to be confirmed with the supervisor. & MALLET, Gibbs sampling. \\
\hline
Iterations & 50 passes over the data & Motivated & Checked against 200 passes for every model: at 50 passes the
  topics match the 200-pass topics with mean cosine 0.996--0.999 and perplexity falls by only 0.1--0.2\% after 50
  passes, so more passes would not change the topics. & 1,000 Gibbs iterations. \\
\hline
Online learning & Mini-batches of 128 motions; learning decay 0.7, offset 10 & Default & scikit-learn defaults; not
  tested. & Not applicable (Gibbs sampling over all documents). \\
\hline
Priors & Symmetric: $1/k$ for both document--topic and topic--word & Default & scikit-learn defaults; not tested. &
  Asymmetric, optimized every 10 iterations. \\
\hline
Random seed & 42 for the reported models & Default & Arbitrary; how much the topics depend on the seed is tested
  separately (section 5). & A fixed random seed. \\
\hline
Shared settings & All models and stability tests take their settings from one module & Motivated & The stability tests
  always measure the models that are reported; a unit test checks that no script fits its own LDA model. &
  Not reported. \\
\hline
%
% ---------------------------------------------------------------- models
\steprow{3 Models}
\hline
Period models & One model per period, each with its own vocabulary: 3,613 / 3,224 / 2,314 / 2,477 motions & Gordon &
  The main analysis: topics are traced across adjacent periods. Periods of complete electoral terms, assigned by
  session. & One model per bin; four bins of 2,143--5,254 chunks. \\
\hline
Whole-corpus model & One model of all 11,628 motions & Gordon & The point of comparison: one set of topics for all
  sessions. Also used for prominence over time, if that analysis is kept. & Full-corpus models with 15 and 60 topics
  as controls. \\
\hline
%
% ---------------------------------------------------------------- topics
\steprow{4 Number of topics}
\hline
Counts tested & 8, 11 and 15 topics per period & Motivated & A small range within the time box of the proposal,
  including Gordon's 15. & 15 per bin, chosen for interpretability and coverage; no range reported. \\
\hline
Selection rule & The highest count at which at least half of the topics are stable (section 5) in every period & Motivated &
  8 topics: 22 of 32 stable (69\%); 11: 24 of 44 (55\%); 15: 28 of 60 (47\%). 11 topics fail the rule in period 1
  (36\%), 15 in every period (47\%). The number of stable topics per period is about the same at every count (4--7), so
  more topics mainly add seed-dependent ones, and the mean cosine is highest at 8 in every period. Every period has at
  least 5 of 8 topics stable. & Not chosen by stability. \\
\hline
Period models & 8 topics per period & Motivated & The result of the selection rule. & 15 per bin. \\
\hline
Whole-corpus model & 32 topics & Motivated & 32 = 4 $\times$ 8: the same total number of topics as the four period
  models together. & 60 = 4 $\times$ 15, the same reasoning; also 15. \\
\hline
%
% ---------------------------------------------------------------- seed stability
\steprow{5 Seed stability (how much the trained topics depend on the random seed)}
\hline
Seeds & Each model refitted with 5 seeds (42, 123, 456, 789, 999) & Motivated & Every model is tested, the period
  models included. & Whole-corpus runs only. \\
\hline
Matching & Topics of two seeds matched one-to-one (Hungarian algorithm) on cosine over the full topic--word
  distributions; overlap of the top 20 words (Jaccard) as a second measure & Motivated & Topic numbers are arbitrary in
  each fitted model, so topics must be matched before they are compared. & Jaccard similarity of the top words across
  seeds. \\
\hline
Stable topic & Found again (matched cosine $\geq 0.70$) in at least 3 of the 4 other seeds' models & Default & A
  judgement call, stated as such; the threshold has not been varied. & No threshold. \\
\hline
Results & Whole-corpus model: 9 of 32 topics stable, mean Jaccard 0.238. Period models: 22 of 32 stable & -- & The
  period models are more stable than the whole-corpus model with the same vocabulary rule; part of the difference is the
  number of topics per model (at 15 per period, 47\% stable). Unstable
  topics are flagged before interpretation. & Whole-corpus runs unstable: mean Jaccard 0.26--0.31. \\
\hline
\end{longtable}
}
\end{landscape}
